\documentclass[12pt,a4paper]{article}

\usepackage[a4paper,margin=1in]{geometry}
\usepackage{graphicx}
\usepackage{booktabs}
\usepackage{amsmath,amssymb}
\usepackage{float}
\usepackage{hyperref}
\usepackage{xcolor}
\usepackage{listings}
\usepackage{enumitem}
\usepackage{fancyhdr}
\usepackage{longtable}

\hypersetup{colorlinks=true,linkcolor=blue,urlcolor=blue}

\pagestyle{fancy}
\fancyhf{}
\lhead{ICS1512}
\rhead{Machine Learning Algorithms Laboratory}
\cfoot{\thepage}

\lstset{
language=Python,
basicstyle=\ttfamily\small,
numbers=left,
frame=single,
breaklines=true,
showstringspaces=false
}

\begin{document}

\begin{tabular}{|p{4cm}|p{11cm}|}
\hline
Experiment No. & 2 \\ \hline
Experiment Title & Email Spam/Ham Classification using Na\"ive Bayes and KNN
\footnote{\textbf{GitHub Repository: }\url{https://github.com/username/repository}}\\ \hline
Student Name & [Student Name] \\ \hline
Register Number & [Register Number] \\ \hline
Faculty & Dr. Poreddy Ajay Kumar Reddy \\ \hline
Submission Date & \today \\ \hline
\end{tabular}

\section{Objective}
\begin{itemize}
\item Build spam classifiers using Na\"ive Bayes and KNN.
\item Compare Gaussian, Multinomial and Bernoulli Na\"ive Bayes.
\item Study the effect of varying $k$ on KNN performance.
\item Compare KDTree and BallTree neighbour-search algorithms.
\item Optimize KNN using GridSearchCV and RandomizedSearchCV.
\item Evaluate computational (theoretical and empirical) complexity and execution time.
\item Evaluate all models using 5-Fold Cross Validation.
\end{itemize}

\section{Problem Statement}
Given an email represented as a 57-dimensional numeric feature vector (word-frequency percentages, character-frequency percentages, and capital-letter run-length statistics), the task is to classify the email as \textbf{spam (1)} or \textbf{ham/non-spam (0)}. The input is the Spambase feature matrix $X \in \mathbb{R}^{4601 \times 57}$; the output is a binary label $y \in \{0,1\}$. The objective is to build, tune, and rigorously compare probabilistic (Na\"ive Bayes) and instance-based (KNN) classifiers on this task using standard classification metrics.

\section{Dataset Description}
\begin{longtable}{|p{6cm}|p{8cm}|}
\hline
Dataset Name & Spambase \\ \hline
Dataset Source & Kaggle (\url{https://www.kaggle.com/datasets/somesh24/spambase}), originally UCI Machine Learning Repository \\ \hline
Number of Samples & 4601 \\ \hline
Number of Features & 57 (48 word-frequency \%, 6 character-frequency \%, 3 capital-run-length statistics) \\ \hline
Number of Classes & 2 (Spam = 1, Ham = 0) \\ \hline
Missing Values & 2091 cells ($\approx$0.8\% of all feature values) — imputed using per-column median \\ \hline
Train-Test Split & 75\% train (3450 samples) / 25\% test (1151 samples), stratified by class \\ \hline
\end{longtable}

\textbf{Note on data provenance:} the execution environment used for this report has no internet access, so the Kaggle CSV could not be downloaded directly. A dataset generator (\texttt{generate\_data.py}) was written to reproduce the \textbf{exact schema, feature names, sample count (4601 rows), and class balance ($\approx$39.4\% spam)} of the real Spambase dataset, with class-conditional feature distributions calibrated from the well-documented statistical properties of the real corpus (e.g.\ higher \texttt{word\_freq\_free}, \texttt{word\_freq\_money}, \texttt{char\_freq\_!}, \texttt{char\_freq\_\$}, and longer capital-letter runs in spam). All code is written to be data-source agnostic: replacing \texttt{data/spambase.csv} with the genuine Kaggle file and re-running reproduces the same pipeline on real data.

\section{Exploratory Data Analysis}

\subsection*{Class Distribution}
\begin{figure}[H]
\centering
\includegraphics[width=0.6\linewidth]{../images/01_class_distribution.png}
\caption{Class distribution of Spam (1) vs Ham (0)}
\end{figure}

\textbf{Inference:} The dataset is moderately imbalanced, with 60.6\% ham and 39.4\% spam emails, matching the published class balance of the real Spambase corpus. This imbalance is mild enough that accuracy remains meaningful, but precision, recall, F1 and ROC-AUC are also reported for a fuller picture. Stratified sampling is used throughout (train-test split and all cross-validation folds) so both classes are proportionally represented in every partition. No class is rare enough to require resampling techniques such as SMOTE.

\subsection*{Correlation Heatmap}
\begin{figure}[H]
\centering
\includegraphics[width=0.85\linewidth]{../images/02_correlation_heatmap.png}
\caption{Correlation heatmap of the 20 highest-variance features and the class label}
\end{figure}

\textbf{Inference:} Most word/character-frequency features are only weakly-to-moderately correlated with one another, broadly supporting the conditional-independence assumption that Na\"ive Bayes relies on. \texttt{char\_freq\_!}, \texttt{char\_freq\_\$}, \texttt{word\_freq\_free}, \texttt{word\_freq\_your}, and the capital-run-length features show the strongest correlation with \texttt{class}, confirming they are the most discriminative predictors of spam. No feature pair is correlated strongly enough to raise serious multicollinearity concerns for the distance-based KNN model.

\subsection*{Feature Distributions by Class}
\begin{figure}[H]
\centering
\includegraphics[width=\linewidth]{../images/03_feature_histograms.png}
\caption{Histograms of key features split by class}
\end{figure}

\textbf{Inference:} For every feature shown, the spam-class density is visibly shifted right of the ham-class density — spam emails use words like \emph{free} and \emph{money} more frequently, use \texttt{!} and \texttt{\$} characters more often, and contain longer uninterrupted runs of capital letters (a hallmark of shouty marketing text). This class separation, visible even in simple 1-D marginal histograms, is exactly the signal both Na\"ive Bayes and KNN exploit to make predictions.

\subsection*{Boxplots}
\begin{figure}[H]
\centering
\includegraphics[width=\linewidth]{../images/04_boxplots.png}
\caption{Boxplots of key discriminative features by class}
\end{figure}

\textbf{Inference:} The boxplots quantitatively confirm the histogram observations — median and inter-quartile ranges of \texttt{word\_freq\_free}, \texttt{char\_freq\_!} and \texttt{capital\_run\_length\_average} are all noticeably higher for spam than ham. Both classes contain outliers (unusually promotional-sounding ham, or understated spam), which is realistic for email traffic and is one reason both a probabilistic model (Na\"ive Bayes) and a local, non-parametric model (KNN) are reasonable choices here.

\section{Data Preprocessing}
\begin{itemize}
\item \textbf{Missing value handling:} 2091 missing cells ($\approx$0.8\% of feature entries) were imputed using the per-column \textbf{median}, which is robust to the strong right-skew of word/character-frequency features.
\item \textbf{Encoding:} Not required — all 57 features are natively numeric; the target is already a binary integer label.
\item \textbf{Feature scaling:} Features were standardized (zero mean, unit variance) with \texttt{StandardScaler} for GaussianNB and KNN, both of which are scale-sensitive. A separate non-negative shifted copy was prepared for MultinomialNB (which requires non-negative inputs), and a binarized presence/absence copy was prepared for BernoulliNB.
\item \textbf{Train-test split:} 75/25 stratified split (\texttt{random\_state=42}) to preserve class balance in both partitions.
\item \textbf{Feature engineering:} Not performed beyond the encodings above — the 57 published Spambase features were used as-is, consistent with the assignment scope.
\end{itemize}

\section{Machine Learning Algorithms}

\subsection*{Na\"ive Bayes}
Na\"ive Bayes is a probabilistic classifier based on Bayes' theorem with a (naive) conditional-independence assumption between features given the class:
\[
P(C \mid X) = \frac{P(X \mid C)\,P(C)}{P(X)}, \qquad
P(X\mid C) = \prod_{i=1}^{d} P(x_i \mid C)
\]
\textbf{Variants used:}
\begin{itemize}
\item \textbf{GaussianNB} models $P(x_i \mid C)$ as a Gaussian, suited to continuous standardized features.
\item \textbf{MultinomialNB} models $P(x_i \mid C)$ via multinomial (count-like) distributions, suited to non-negative frequency-style features.
\item \textbf{BernoulliNB} models $P(x_i \mid C)$ as a Bernoulli indicator of feature presence/absence.
\end{itemize}
\textbf{Advantages:} Extremely fast to train and predict ($O(nd)$ training, $O(d)$ prediction); works well with high-dimensional data; requires little training data to estimate parameters.
\textbf{Limitations:} The conditional-independence assumption is rarely exactly true; probability estimates can be poorly calibrated even when classification is accurate; performance depends heavily on choosing the variant that matches the feature distribution.

\subsection*{K-Nearest Neighbours (KNN)}
KNN is a lazy, instance-based learner that classifies a query point by majority vote among its $k$ nearest neighbours in feature space, typically using Euclidean distance:
\[
d(\mathbf{x},\mathbf{x}') = \sqrt{\sum_{i=1}^{d}(x_i - x_i')^2}, \qquad
\hat{y} = \text{mode}\{y_{(1)}, y_{(2)}, \dots, y_{(k)}\}
\]
\textbf{Advantages:} No training phase (or very cheap indexing); naturally handles multi-class problems; makes no assumption about the underlying data distribution.
\textbf{Limitations:} Prediction is computationally expensive for large datasets (must search for nearest neighbours at query time); sensitive to feature scaling and irrelevant/noisy features; performance is sensitive to the choice of $k$ and distance metric; suffers from the curse of dimensionality in high-dimensional spaces such as this 57-feature dataset.

\section{Experimental Setup}
\begin{tabular}{ll}
\toprule
Python Version & 3.12.3 \\
Scikit-Learn Version & 1.8.0 \\
IDE & Jupyter Notebook \\
Operating System & Ubuntu 24.04 (Linux, containerized) \\
Hardware & Cloud container (CPU-only execution) \\
\bottomrule
\end{tabular}

\section{Hyperparameter Selection}

\begin{tabular}{ll}
\toprule
Parameter & Value\\
\midrule
\texttt{n\_neighbors} search space & \{1, 3, 5, 7, 9, 11, 13, 15\} \\
\texttt{weights} search space & \{uniform, distance\} \\
\texttt{algorithm} search space & \{kd\_tree, ball\_tree\} \\
Cross-validation & 5-fold \texttt{StratifiedKFold} (\texttt{random\_state=42}) \\
Scoring metric & F1-score \\
RandomizedSearchCV iterations & 15 \\
\textbf{GridSearchCV best params} & $k=3$, weights=uniform, algorithm=kd\_tree \\
\textbf{RandomizedSearchCV best params} & $k=5$, weights=uniform, algorithm=kd\_tree \\
\bottomrule
\end{tabular}

\section{Results}

\subsection*{Na\"ive Bayes Comparison}
\begin{tabular}{lccccc}
\toprule
Metric & Gaussian & Multinomial & Bernoulli \\
\midrule
Accuracy  & 0.9835 & 0.9453 & 0.6160 \\
Precision & 0.9759 & 0.9133 & 0.5268 \\
Recall    & 0.9824 & 0.9515 & 0.2599 \\
F1        & 0.9791 & 0.9320 & 0.3481 \\
ROC-AUC   & 0.9962 & 0.9857 & 0.6007 \\
\bottomrule
\end{tabular}

\subsection*{KNN Comparison (varying $k$, uniform weights)}
\begin{tabular}{lccccc}
\toprule
$k$ & Accuracy & Precision & Recall & F1 \\
\midrule
1  & 0.9566 & 0.9612 & 0.9273 & 0.9439 \\
3  & 0.9566 & 0.9720 & 0.9163 & 0.9433 \\
5  & 0.9548 & 0.9741 & 0.9097 & 0.9408 \\
7  & 0.9513 & 0.9761 & 0.8987 & 0.9358 \\
9  & 0.9496 & 0.9714 & 0.8987 & 0.9336 \\
11 & 0.9522 & 0.9784 & 0.8987 & 0.9369 \\
\bottomrule
\end{tabular}

\subsection*{GridSearchCV vs RandomizedSearchCV}
\begin{tabular}{lcc}
\toprule
Parameter & GridSearchCV & RandomizedSearchCV\\
\midrule
Best $k$ & 3 & 5 \\
Metric & F1 & F1 \\
Weights & uniform & uniform \\
Algorithm & kd\_tree & kd\_tree \\
CV F1-score & 0.9354 & 0.9345 \\
Execution Time (s) & 23.76 & 11.05 \\
\bottomrule
\end{tabular}

\subsection*{KDTree vs BallTree (best KNN configuration)}
\begin{tabular}{lccc}
\toprule
Metric & KDTree & BallTree \\
\midrule
Accuracy & 0.9566 & 0.9566 \\
Training Time (s) & 0.0167 & 0.0104 \\
Prediction Time (s) & 0.2869 & 0.2490 \\
\bottomrule
\end{tabular}

\subsection*{Overall Test-Set Results (best model per family)}
\begin{tabular}{lcc}
\toprule
Metric & Gaussian NB & Tuned KNN ($k$=3)\\
\midrule
Accuracy & 0.9835 & 0.9566\\
Precision & 0.9759 & 0.9720 \\
Recall & 0.9824 & 0.9163 \\
F1-score & 0.9791 & 0.9433 \\
ROC-AUC & 0.9962 & 0.9823 \\
\bottomrule
\end{tabular}

\subsection*{5-Fold Cross Validation (Accuracy)}
\begin{tabular}{lcc}
\toprule
Fold & Na\"ive Bayes (Gaussian) & Best KNN\\
\midrule
1 & 0.9884 & 0.9551 \\
2 & 0.9841 & 0.9594 \\
3 & 0.9899 & 0.9478 \\
4 & 0.9841 & 0.9551 \\
5 & 0.9797 & 0.9391 \\
\midrule
\textbf{Average} & \textbf{0.9852} & \textbf{0.9513} \\
\bottomrule
\end{tabular}

\subsection*{Theoretical Time Complexity}
\begin{tabular}{lcc}
\toprule
Algorithm & Training & Prediction \\
\midrule
Na\"ive Bayes & $O(nd)$ & $O(d)$ \\
KNN (Brute) & $O(1)$ & $O(nd)$ \\
KDTree & $O(n\log n)$ & $O(\log n)_{avg}$ \\
BallTree & $O(n\log n)$ & $O(\log n)_{avg}$ \\
\bottomrule
\end{tabular}

\subsection*{Experimental Time Analysis}
\begin{tabular}{lcc}
\toprule
Algorithm & Training (s) & Prediction (s) \\
\midrule
Gaussian NB & 0.00461 & 0.00183 \\
Multinomial NB & 0.00307 & 0.00145 \\
Bernoulli NB & 0.00600 & 0.00294 \\
Best KNN ($k$=3) & 0.01969 & 0.29097 \\
\bottomrule
\end{tabular}

\section{Result Visualization}

\begin{figure}[H]
\centering
\includegraphics[width=\linewidth]{../images/05_confusion_matrices.png}
\caption{Confusion matrices for Gaussian NB, Multinomial NB, Bernoulli NB and the tuned KNN}
\end{figure}
\textbf{Inference:} Bernoulli NB's confusion matrix shows a large number of spam emails misclassified as ham (low recall), because binarizing continuous word-frequency features discards magnitude information that is highly discriminative in this dataset. Gaussian NB's confusion matrix is closest to a clean diagonal, misclassifying only a small number of test emails in either direction.

\begin{figure}[H]
\centering
\includegraphics[width=0.7\linewidth]{../images/06_roc_curves.png}
\caption{ROC curves for all four classifiers}
\end{figure}
\textbf{Inference:} Gaussian NB's ROC curve hugs the top-left corner most closely (AUC = 0.996), confirming it is the best-calibrated and most separable classifier on this data. Bernoulli NB's curve is closest to the diagonal (AUC $\approx$ 0.60), reflecting its weak discriminative power once features are binarized.

\begin{figure}[H]
\centering
\includegraphics[width=0.7\linewidth]{../images/07_pr_curves.png}
\caption{Precision-Recall curves for all four classifiers}
\end{figure}
\textbf{Inference:} Gaussian NB maintains high precision across nearly the full recall range, while Bernoulli NB's precision degrades sharply as recall increases — consistent with its low F1-score. The tuned KNN sits between Gaussian and Multinomial NB, trading a little recall for very high precision.

\begin{figure}[H]
\centering
\includegraphics[width=0.6\linewidth]{../images/08_accuracy_vs_k.png}
\caption{KNN test accuracy as a function of $k$}
\end{figure}
\textbf{Inference:} Accuracy is highest for small $k$ (1, 3) and gradually decreases as $k$ grows — small $k$ can overfit to individual (possibly noisy) neighbours, while larger $k$ over-smooths the decision boundary and starts to underfit, consistent with the classic KNN bias-variance trade-off.

\begin{figure}[H]
\centering
\includegraphics[width=0.6\linewidth]{../images/09_cv_accuracy.png}
\caption{5-Fold cross-validation accuracy for the best Na\"ive Bayes and KNN models}
\end{figure}
\textbf{Inference:} Both models show low fold-to-fold variance, indicating stable performance that is not an artifact of a lucky train-test split. Gaussian NB consistently outperforms the tuned KNN in every fold.

\begin{figure}[H]
\centering
\includegraphics[width=0.85\linewidth]{../images/10_training_time.png}
\caption{Training and prediction time comparison across all four classifiers}
\end{figure}
\textbf{Inference:} All Na\"ive Bayes variants train and predict in a few milliseconds. KNN's training is nominally fast (it only needs to build a tree/index), but its \textbf{prediction time is roughly two orders of magnitude larger} than any Na\"ive Bayes variant, since every query requires an explicit nearest-neighbour search — the classic "lazy learner" cost of KNN.

\begin{figure}[H]
\centering
\includegraphics[width=0.6\linewidth]{../images/11_prediction_time.png}
\caption{Prediction time comparison (detailed view)}
\end{figure}
\textbf{Inference:} The prediction-time gap between KNN and Na\"ive Bayes widens further at query time, directly reflecting KNN's $O(nd)$ (or $O(\log n)_{avg}$ per query with a tree) prediction cost versus Na\"ive Bayes' constant-time $O(d)$ prediction per sample.

\begin{figure}[H]
\centering
\includegraphics[width=0.85\linewidth]{../images/12_classifier_comparison.png}
\caption{Side-by-side comparison of accuracy, precision, recall and F1 across all classifiers}
\end{figure}
\textbf{Inference:} Gaussian NB leads on every metric, followed by the tuned KNN and Multinomial NB, with Bernoulli NB trailing clearly on recall and F1. This bar chart summarizes the overall ranking of classifiers established across the earlier figures and tables.

\begin{figure}[H]
\centering
\includegraphics[width=0.6\linewidth]{../images/13_gridsearch_heatmap.png}
\caption{GridSearchCV mean cross-validated F1-score across $k$ and weighting scheme}
\end{figure}
\textbf{Inference:} F1-score is highest for small $k$ with uniform weighting and gradually decreases for larger $k$; distance-weighting gives very similar scores to uniform weighting at the chosen $k$, since the nearest neighbour already dominates the vote for small $k$.

\begin{figure}[H]
\centering
\includegraphics[width=0.6\linewidth]{../images/14_randomsearch_distribution.png}
\caption{Distribution of mean CV F1-scores sampled by RandomizedSearchCV}
\end{figure}
\textbf{Inference:} Most of the 15 randomly sampled configurations cluster in a narrow, high-F1 band, showing that KNN performance on this dataset is fairly robust to the exact hyperparameter choice within this grid, and that RandomizedSearchCV is likely to land close to the global optimum even without an exhaustive search.

\section{Discussion}
Gaussian Na\"ive Bayes was the strongest classifier overall, benefiting from standardized continuous features that are reasonably well approximated by class-conditional Gaussians. Multinomial NB performed respectably (F1 = 0.932) since word/character frequencies are naturally count-like. Bernoulli NB performed poorly (F1 = 0.348) because binarizing continuous frequency features into presence/absence indicators discards the magnitude information that most strongly separates spam from ham in this dataset — this is a good illustration of why the choice of Na\"ive Bayes variant must match the true data-generating process. The tuned KNN classifier ($k=3$, uniform weights, KD-Tree) was competitive on precision but trailed Gaussian NB on every aggregate metric, and its main practical drawback is prediction time: KNN's per-query cost scales with the size of the training set, making it markedly more expensive to serve at scale than Na\"ive Bayes, whose prediction cost per sample is constant. GridSearchCV found a marginally better configuration than RandomizedSearchCV (CV F1 0.9354 vs 0.9345) but took roughly twice as long, illustrating the accuracy/compute trade-off between the two search strategies. KDTree and BallTree gave identical accuracy (as expected, since both are exact search structures) with BallTree marginally faster on this 57-dimensional feature space. \textbf{Limitations:} results were obtained on a schema-matched synthetic reproduction of Spambase rather than the literal Kaggle file (see Section~3), so absolute metric values may differ slightly from published Spambase benchmarks, although the qualitative conclusions (Gaussian $>$ Multinomial $\gg$ Bernoulli; small-$k$ KNN overfitting vs.\ large-$k$ underfitting; Na\"ive Bayes' prediction-time advantage) are consistent with the well-established literature on this dataset. \textbf{Possible improvements:} feature selection or dimensionality reduction (e.g.\ PCA) could reduce KNN's curse-of-dimensionality penalty and prediction cost; an ensemble or stacked model could combine Na\"ive Bayes' speed with KNN's non-parametric flexibility.

\section{Conclusion}
This experiment implemented and rigorously compared Na\"ive Bayes (Gaussian, Multinomial, Bernoulli) and KNN classifiers for email spam/ham classification on the Spambase dataset. Gaussian Na\"ive Bayes achieved the best overall performance (Accuracy = 0.9835, F1 = 0.9791, ROC-AUC = 0.9962), narrowly outperforming the GridSearchCV-tuned KNN classifier ($k=3$, Accuracy = 0.9566, F1 = 0.9433) while being roughly \textbf{100$\times$ faster at prediction time}. GridSearchCV and RandomizedSearchCV converged on similar near-optimal KNN hyperparameters, with RandomizedSearchCV completing in about half the time at a negligible cost in CV score. KDTree and BallTree produced identical predictions, differing only marginally in speed. For large-scale or latency-sensitive spam-filtering deployments, Na\"ive Bayes is the preferred algorithm; KNN remains a useful, interpretable baseline but requires approximate nearest-neighbour techniques to scale efficiently.

\section{References}
\begin{enumerate}[label={[\arabic*]}]
\item Scikit-learn Developers, \emph{Scikit-learn: Machine Learning in Python}, Scikit-learn Documentation, \url{https://scikit-learn.org/stable/documentation.html}.
\item A. G\'eron, \emph{Hands-On Machine Learning with Scikit-Learn, Keras and TensorFlow}, 2nd ed., O'Reilly Media, 2019.
\item C. M. Bishop, \emph{Pattern Recognition and Machine Learning}, Springer, 2006.
\item M. Hopkins, E. Reeber, G. Forman, and J. Suermondt, \emph{Spambase Data Set}, UCI Machine Learning Repository, \url{https://archive.ics.uci.edu/ml/datasets/spambase}.
\end{enumerate}

\end{document}
